\rtmtight
\begin{tblr}{width=\textwidth,colspec={|Q[c,m,2.1cm]|X[l,m]|},hlines,vlines,hspan=minimal,rowsep=1.5pt,colsep=3pt,row{1}={bg=headgray,font=\bfseries}}
\SetCell{c,m}\hfil\logo\hfil & \textbf{POLITEKNIK NEGERI MALANG}\par\textbf{JURUSAN TEKNOLOGI INFORMASI}\par\textbf{PROGRAM STUDI: D4 TEKNIK INFORMATIKA} \\
\SetCell[c=2]{l} \textbf{RENCANA TUGAS MAHASISWA (RTM) DAN RUBRIK PENILAIAN} & \\
Nama Mata Kuliah & Praktikum Dasar Pemrograman \\
Kode Mata Kuliah & RTI251007 \quad \textbf{sks / Jam perminggu:} 3 SKS/6 jam \quad \textbf{Semester:} 1 \\
Dosen Pengampu & \begin{enumerate}\item Mungki Astiningrum, S.T., M.Kom.\item Imam Fahrur Rozi, S.T., M.T.\item Triana Fatmawati, S.T., M.T.\item Mamluatul Hani'ah, S.Kom., M.Kom.\item Rudy Ariyanto, S.T., M.Cs.\item Noprianto, S.Kom., M.Eng.\end{enumerate} \\
\end{tblr}
\assessmentblock{Tugas Praktikum Jobsheet Dasar Pemrograman}{Tugas praktikum (jobsheet)}{SCPMK0903-00701, SCPMK0704-00702, SCPMK0704-00703, SCPMK0704-00704}{Mahasiswa menyelesaikan jobsheet praktikum mingguan (instalasi Java, studi kasus algoritma, sequence, pemilihan, perulangan, array, fungsi) beserta tugas mandiri tiap pertemuan.}{Melakukan review materi bersama dosen, melaksanakan langkah percobaan sesuai jobsheet, melakukan compile dan debugging program Java, mengerjakan tugas praktikum, dan mengumpulkan laporan melalui LMS.}{Laporan jobsheet praktikum (PDF), kode sumber Java, dan screenshot output program yang dikumpulkan melalui LMS.}{Ketepatan langkah percobaan dan jawaban pertanyaan jobsheet & 8 & Seluruh langkah percobaan dijalankan sesuai jobsheet dan pertanyaan jobsheet dijawab benar disertai penjelasan. & Sebagian besar langkah percobaan dijalankan dan sebagian pertanyaan jobsheet dijawab benar. & Langkah percobaan tidak diselesaikan atau jawaban pertanyaan jobsheet banyak yang salah. \\ \\
Ketepatan kode program dan output tugas praktikum & 10 & Kode program dapat di-compile, berjalan tanpa error, dan output sesuai spesifikasi seluruh tugas praktikum. & Kode program berjalan tetapi sebagian output belum sesuai spesifikasi atau masih terdapat kesalahan minor. & Kode program tidak dapat di-compile, tidak berjalan, atau output tidak sesuai spesifikasi tugas. \\ \\
Kedisiplinan pengumpulan dan keaktifan diskusi & 4 & Seluruh laporan dikumpulkan tepat waktu dan mahasiswa aktif dalam diskusi/tanya jawab daring maupun luring. & Sebagian laporan terlambat dikumpulkan atau keaktifan diskusi terbatas. & Laporan banyak yang tidak dikumpulkan dan mahasiswa pasif dalam diskusi. \\}

\assessmentblock{Kuis 1 Tes Praktik Dasar Pemrograman}{Kuis (tes praktik)}{SCPMK0903-00701, SCPMK0704-00702}{Mahasiswa menyelesaikan tes praktik individu materi minggu ke-1 sampai 3: pemodelan algoritma, tipe data, variabel, operator, input-output, dan sequence.}{Menganalisis soal studi kasus, memodelkan penyelesaian dalam algoritma/flowchart, menuliskan program Java, serta melakukan compile dan uji output dalam waktu yang ditentukan.}{Berkas algoritma/flowchart dan kode sumber Java beserta output program yang dikumpulkan melalui LMS.}{Ketepatan pemodelan algoritma (input, proses, output) & 4 & Algoritma/flowchart menggambarkan input, proses, dan output secara runtut dan sesuai studi kasus. & Algoritma/flowchart menggambarkan penyelesaian tetapi terdapat langkah yang kurang tepat atau tidak runtut. & Algoritma/flowchart tidak menggambarkan penyelesaian studi kasus. \\ \\
Ketepatan implementasi tipe data, operator, input-output, dan sequence & 6 & Program Java menggunakan tipe data, variabel, operator, dan input-output secara tepat; kode berjalan dan output benar. & Program berjalan tetapi terdapat pemilihan tipe data/operator yang kurang tepat atau output sebagian salah. & Program tidak dapat di-compile atau output tidak sesuai soal. \\}

\assessmentblock{UTS Ujian Praktik Pemilihan dan Perulangan}{UTS (ujian praktik)}{SCPMK0704-00702, SCPMK0704-00703}{Mahasiswa menyelesaikan ujian praktik individu pembuatan program Java yang memuat struktur pemilihan dan perulangan berdasarkan studi kasus.}{Menganalisis studi kasus ujian, merancang penyelesaian, mengimplementasikan program Java dengan struktur pemilihan dan perulangan, serta menguji kebenaran output dalam waktu ujian.}{Kode sumber Java dan hasil eksekusi program yang dikumpulkan melalui LMS pada akhir sesi ujian.}{Ketepatan analisis kasus dan rancangan penyelesaian & 6 & Input, proses, dan output kasus diidentifikasi tepat dan rancangan penyelesaian sesuai kebutuhan soal. & Analisis kasus sebagian tepat tetapi terdapat kebutuhan soal yang tidak terjawab. & Analisis kasus salah sehingga penyelesaian tidak menjawab soal. \\ \\
Ketepatan implementasi struktur pemilihan & 7 & Struktur pemilihan (if, if-else, else-if, switch-case, bersarang) dipilih dan diimplementasikan tepat sesuai kondisi kasus. & Struktur pemilihan diimplementasikan tetapi terdapat kondisi yang salah atau tidak lengkap. & Struktur pemilihan salah total atau tidak digunakan padahal dibutuhkan. \\ \\
Ketepatan implementasi struktur perulangan dan output program & 7 & Struktur perulangan diimplementasikan tepat; program berjalan tanpa error dan seluruh output sesuai soal. & Perulangan berjalan tetapi batas/kondisi kurang tepat sehingga sebagian output salah. & Program tidak berjalan atau output tidak sesuai soal. \\}

\assessmentblock{Kuis 2 Tes Praktik Perulangan Bersarang dan Array}{Kuis (tes praktik)}{SCPMK0704-00703}{Mahasiswa menyelesaikan tes praktik individu materi perulangan bersarang, array 1 dimensi, dan array 2 dimensi.}{Menganalisis soal studi kasus, mengimplementasikan program Java dengan perulangan bersarang dan array, serta menguji kebenaran output dalam waktu yang ditentukan.}{Kode sumber Java beserta output program yang dikumpulkan melalui LMS.}{Ketepatan implementasi perulangan bersarang & 5 & Perulangan bersarang diimplementasikan tepat, kode berjalan, dan output sesuai studi kasus. & Perulangan bersarang berjalan tetapi batas/kondisi kurang tepat sehingga output sebagian salah. & Perulangan bersarang tidak diimplementasikan atau program tidak berjalan. \\ \\
Ketepatan implementasi array 1 dimensi dan 2 dimensi & 5 & Deklarasi, pengisian, dan pengaksesan elemen array tepat; output pengolahan data sesuai soal. & Array digunakan tetapi terdapat kesalahan pengaksesan elemen atau output sebagian salah. & Array tidak digunakan dengan benar atau output tidak sesuai soal. \\}

\assessmentblock{PBL Proyek Akhir Pemrograman}{Project Based Learning}{SCPMK0704-00704}{Mahasiswa membangun program Java utuh untuk menyelesaikan studi kasus permasalahan sesuai topik yang ditentukan dengan memanfaatkan materi pertemuan 1--14.}{Menganalisis permasalahan, membuat flowchart penyelesaian, mengimplementasikan program Java (sequence, pemilihan, perulangan, array, fungsi), menguji program, dan menyiapkan demo.}{Flowchart, kode sumber Java, dan dokumentasi singkat proyek yang dikumpulkan melalui LMS.}{Ketepatan rancangan flowchart proyek & 3 & Flowchart menggambarkan alur penyelesaian studi kasus secara runtut, lengkap, dan konsisten dengan program. & Flowchart tersedia tetapi terdapat alur yang kurang lengkap atau tidak konsisten dengan program. & Flowchart tidak tersedia atau tidak menggambarkan penyelesaian studi kasus. \\ \\
Ketepatan implementasi program proyek & 3 & Program memanfaatkan pemilihan, perulangan, array, dan fungsi secara tepat; kode berjalan sesuai studi kasus. & Program berjalan tetapi pemanfaatan struktur pemrograman kurang tepat atau sebagian fitur belum selesai. & Program tidak berjalan atau tidak sesuai studi kasus. \\ \\
Kelengkapan dokumentasi dan kesiapan demo & 2 & Dokumentasi proyek lengkap dan mahasiswa siap mendemonstrasikan program dengan lancar. & Dokumentasi atau kesiapan demo masih terbatas. & Tidak ada dokumentasi dan mahasiswa tidak siap demo. \\}

\assessmentblock{UAS Demo Project Dasar Pemrograman}{UAS (ujian praktik dan demo project)}{SCPMK0704-00704}{Mahasiswa mendemonstrasikan program proyek akhir di hadapan penguji dan menjawab pertanyaan terkait logika serta struktur kode program.}{Menjalankan program di hadapan penguji, menunjukkan fungsionalitas sesuai studi kasus, menjelaskan alur logika dan struktur kode, serta menjawab pertanyaan modifikasi sederhana.}{Demo program berjalan, kode sumber Java, dan jawaban lisan atas pertanyaan penguji.}{Kelengkapan fungsionalitas program & 12 & Seluruh fungsionalitas sesuai studi kasus berjalan benar tanpa error saat didemonstrasikan. & Fungsionalitas utama berjalan tetapi sebagian fitur error atau belum selesai. & Program tidak berjalan atau fungsionalitas utama tidak sesuai studi kasus. \\ \\
Kualitas kode program (struktur, fungsi, kerapian) & 10 & Kode terstruktur dengan fungsi yang tepat, penamaan konsisten, dan mudah dibaca. & Kode berjalan tetapi struktur fungsi atau kerapian penulisan masih kurang. & Kode tidak terstruktur, sulit dibaca, atau tidak memanfaatkan fungsi. \\ \\
Kualitas demo dan penjelasan program & 8 & Mahasiswa mendemonstrasikan program dengan lancar dan mampu menjelaskan serta memodifikasi kode sesuai pertanyaan penguji. & Demo berjalan tetapi penjelasan logika kode kurang meyakinkan atau modifikasi tidak selesai. & Mahasiswa tidak mampu menjelaskan kode program yang didemonstrasikan. \\}

\begin{tblr}{width=\textwidth,colspec={|Q[l,m,3.2cm]|X[l,m]|},hlines,vlines,hspan=minimal,row{1-3}={bg=headgray},rowsep=1.5pt,colsep=3pt}
Jadwal Pelaksanaan: & Minggu 1--17 sesuai RPS. Setiap evaluasi memiliki RTM dan rubrik multi-kriteria. \\
Lain-Lain Yang Diperlukan: & LMS (lms.polinema.ac.id), JDK dan IDE, komputer laboratorium, dan perangkat presentasi. \\
Pustaka: & Mengacu pustaka utama dan pendukung pada bagian RPS. \\
\end{tblr}
